% =====================================================================
%  3 Geschäftsregeln  (Aufgabenstellung: „Beschreibung der wichtigsten
%  Regeln für Geschäftsobjekte und Geschäftsprozesse, die zwingend
%  eingehalten werden müssen (z. B., dass das Startdatum eines Vertrags
%  immer vor dessen Enddatum liegen muss)“)
%  Jede Regel später im Testdokument prüfen (GR-ID im Testfall nennen).
% =====================================================================
\section{Geschäftsregeln}\label{sec:geschaeftsregeln}
Geschäftsregeln sind fachliche Bedingungen, die BalanceOS unabhängig vom Weg der Eingabe stets einhält. Tabelle~\ref{tab:geschaeftsregeln} fasst sie nach Geschäftsobjekten gruppiert zusammen. Jede Regel ist dort den betroffenen Objekten bzw. Prozessen aus den Abschnitten~\ref{sec:geschaeftsobjekte} und~\ref{sec:prozessuebersicht} sowie den User Stories zugeordnet. Verstößt eine Eingabe gegen eine Regel, wird sie vor dem Speichern abgewiesen und am betroffenen Feld erläutert (NFA02). Die Gültigkeitsregeln der Eingabefelder und die Fehlermeldungen dazu enthält Tabelle~\ref{tab:validierung} in Abschnitt~\ref{sec:validierung}. Im Testdokument prüft mindestens ein Testfall jede Geschäftsregel und verweist dabei auf ihre Kennung.

\begin{longtable}{@{}L{1.2cm}L{9.0cm}L{3.5cm}L{1.5cm}@{}}
  \caption{Geschäftsregeln}\label{tab:geschaeftsregeln}\\
  \toprule
  \textbf{ID} & \textbf{Regel} & \textbf{Objekte und Prozesse} & \textbf{US} \\
  \midrule
  \endfirsthead
  \toprule
  \textbf{ID} & \textbf{Regel} & \textbf{Objekte und Prozesse} & \textbf{US} \\
  \midrule
  \endhead
  \bottomrule
  \endfoot
  \multicolumn{4}{@{}l}{\textbf{Lebensbereiche und Wunschprofil}} \\*
  GR01 & Der Name eines Lebensbereichs ist 1--40 Zeichen lang (ohne führende und abschließende Leerzeichen) und eindeutig; Groß- und Kleinschreibung werden beim Vergleich nicht unterschieden. & Lebensbereich & US02 \\
  GR02 & Es gibt stets mindestens einen aktiven Lebensbereich mit einem Soll größer als 0~h. Eine Änderung, die diese Bedingung verletzt, wird abgewiesen. & Lebensbereich, GP03 & US02, US04 \\
  GR03 & Die zehn Standardbereiche (istStandard) können umbenannt, umgefärbt, im Typ geändert und deaktiviert, aber nicht gelöscht werden. & Lebensbereich & US01, US02, US03 \\
  GR04 & Ein eigener Lebensbereich, dem Aktivitäten oder Wiederholungsserien zugeordnet sind, wird nicht gelöscht, sondern nur deaktiviert. Löschen ist nur ohne Zuordnungen möglich. & Lebensbereich, Aktivitaet & US02 \\
  GR05 & Aktivitäten deaktivierter Lebensbereiche bleiben erhalten, werden aber nicht bewertet. Neue Aktivitäten können nur aktiven Bereichen zugeordnet werden. & Lebensbereich, GP01, GP03 & US02, US05 \\
  GR06 & Das Soll je Lebensbereich liegt zwischen 0 und 168~h pro Woche; die Summe der Soll-Werte aller aktiven Bereiche überschreitet 168~h nicht. & Lebensbereich & US04 \\
  GR07 & Das Gewicht eines Lebensbereichs ist eine ganze Zahl von 0 bis 10 (Standard~1). & Lebensbereich, GP05 & US19 \\
  GR08 & Der Parameter $\alpha$ liegt zwischen 0 und 1 (Standard 0,6). & Einstellungen, GP03 & US11 \\
  \midrule
  \multicolumn{4}{@{}l}{\textbf{Aktivitäten}} \\*
  GR09 & Jede Aktivität gehört zu genau einem Lebensbereich; eine Mehrfachzuordnung ist nicht möglich. Der Titel umfasst 1--80 Zeichen, die Notiz höchstens 500 Zeichen. & Aktivitaet, GP01 & US05 \\
  GR10 & Die Dauer beträgt 1--1.440~min. Sind Beginn und Ende angegeben, wird die Dauer daraus berechnet. Ein Ende ohne Beginn ist unzulässig, und Beginn und Ende dürfen nicht gleich sein. & Aktivitaet, GP01 & US05 \\
  GR11 & Liegt das Ende vor dem Beginn, verläuft die Aktivität über Mitternacht: Die Dauer ist Ende $-$ Beginn $+$ 24~h, die Aktivität zählt vollständig zum Datum des Beginns (z.\,B. Schlaf 23:00--06:30). & Aktivitaet, GP01, GP03 & US05 \\
  GR12 & Der Status DURCHGEFUEHRT ist für Daten in der Zukunft nicht zulässig. & Aktivitaet, GP01 & US05, US06 \\
  GR13 & Zufriedenheit, Belastung und Erholung sind ganze Zahlen von 1 bis 5 und nur bei Status DURCHGEFUEHRT erfassbar. Wechselt der Status, werden vorhandene Bewertungen entfernt. & Aktivitaet & US05, US06 \\
  GR14 & In Wochenbewertung, Vollständigkeit und Regelprüfung zählen ausschließlich Aktivitäten mit Status DURCHGEFUEHRT in aktiven Lebensbereichen. GEPLANT und AUSGEFALLEN erscheinen nur in der Wochen- und Aktivitätenansicht. & Aktivitaet, GP03, GP04 & US09, US11, US12 \\
  GR15 & Eine Bewegungsart kann nur bei Aktivitäten im Standardbereich Bewegung angegeben werden. R01 und R02 berücksichtigen nur Aktivitäten mit Bewegungsart. & Aktivitaet, GP04 & US05, US12 \\
  \midrule
  \multicolumn{4}{@{}l}{\textbf{Wiederholungsserien}} \\*
  GR16 & Eine Serie hat mindestens einen Wochentag und eine Dauer von 1--1.440~min. Das Enddatum ist optional; ist es angegeben, liegt es nicht vor dem Startdatum. & Wiederholungs\-serie & US07 \\
  GR17 & Termine entstehen nur an den gewählten Wochentagen ab dem Startdatum, höchstens bis zum Enddatum und höchstens 52~Wochen im Voraus. & Wiederholungs\-serie, GP02 & US07 \\
  GR18 & Eine Ausnahme ist nur für ein Termindatum der Serie zulässig (Datum im Gültigkeitszeitraum und an einem gewählten Wochentag), je Datum höchstens einmal. Für ein Datum mit Ausnahme entsteht kein Termin. Auch das Löschen oder Verschieben eines einzelnen Termins legt für sein ursprüngliches Datum eine Ausnahme an. & Ausnahme, GP02 & US08 \\
  GR19 & Das Bearbeiten eines einzelnen Termins ändert nur diesen Termin. Änderungen an der Serie wirken nur auf zukünftige, nicht einzeln bearbeitete Termine im Status GEPLANT. & Wiederholungs\-serie, Aktivitaet & US06, US07 \\
  GR20 & Beim Löschen einer Serie werden ihre zukünftigen, nicht einzeln bearbeiteten Termine im Status GEPLANT und alle Ausnahmen entfernt; übrige Termine bleiben als einzelne Aktivitäten erhalten. & Wiederholungs\-serie, Ausnahme & US06, US07 \\
  \midrule
  \multicolumn{4}{@{}l}{\textbf{Bewertung und Hinweise}} \\*
  GR21 & Bei einer Vollständigkeit unter 4 von 7 Tagen werden weder Bereichs- noch Gesamt-Scores angezeigt; es erscheint nur der Systemhinweis R10. & Wochen\-bewertung, GP03 & US10, US11 \\
  GR22 & Lebensbereiche mit einem Soll von 0~h werden nicht bewertet und gehen nicht in den Gesamt-Score ein. & Bereichs\-bewertung, GP03 & US11 \\
  GR23 & Der Gesamt-Score wird nur bei aktivierter Einstellung gesamtScoreAnzeigen und nur über bewertete Bereiche mit Gewicht größer 0 berechnet. Er ist stets als \enquote{programmeigener Score, nicht validiert} gekennzeichnet. & Wochen\-bewertung, GP05 & US19 \\
  GR24 & Je Regel und Kalenderwoche entsteht höchstens ein Hinweis, bei bereichsbezogenen Regeln höchstens einer je Lebensbereich. Ein ausgeblendeter Hinweis erscheint für dieselbe Woche nicht erneut. & Hinweis, GP04 & US12 \\
  GR25 & Deaktivierte Regeln erzeugen keine Hinweise. Jeder Hinweis enthält eine Begründung mit Regel, Datenbasis, Zeitraum, Schwelle, Quelle und Evidenzart und stellt keine Diagnose dar. & Regel, Hinweis, GP04 & US13, US14 \\
  \midrule
  \multicolumn{4}{@{}l}{\textbf{Fragebogen, Reflexion und Datenverwaltung}} \\*
  GR26 & Eine WHO-5-Erhebung wird nur mit allen fünf Antworten (je 0--5) gespeichert; ihr Datum liegt nicht in der Zukunft. & Who5Erhebung, GP06 & US20 \\
  GR27 & Je Kalenderwoche gibt es höchstens eine Wochenreflexion. & Wochen\-reflexion & US22 \\
  GR28 & Demodaten werden nur zusätzlich zu vorhandenen Daten geladen und verändern weder eigene Aktivitäten noch Wunschprofil oder Einstellungen. \enquote{Demodaten entfernen} löscht ausschließlich Aktivitäten mit istDemo = wahr. Es ist höchstens ein Demodatensatz gleichzeitig geladen. & Aktivitaet, GP09 & US18 \\
  GR29 & \enquote{Alle Daten löschen} wird erst nach Bestätigung durch Eingabe des Wortes LÖSCHEN ausgeführt. Danach sind keine Nutzdaten mehr vorhanden; Standardbereiche, Regeln und Einstellungen entsprechen dem ersten Start. & alle Objekte, GP08 & US17 \\
  GR30 & Daten verlassen den Rechner nur durch einen ausdrücklich ausgelösten Export in eine von Nutzer:innen gewählte Datei; es findet kein Netzwerkzugriff statt. & GP07 & US16 \\
\end{longtable}
\quelle{Eigene Darstellung.}
